\chapter{Applied \& Computational Mathematics}

This document provides a foundational overview of the topics in the Applied \& Computational Mathematics track. For each subsection, essential prerequisite knowledge, core definitions, fundamental theorems, and illustrative examples are presented to guide beginner students.

\section{Calculus}
\textbf{Prerequisite Knowledge:} High school algebra, trigonometry, coordinate geometry, and the concept of limits.

\textbf{Definitions:}
\begin{itemize}
    \item \textbf{Derivative:} The rate of change of a function with respect to a variable, defined as $f'(x) = \lim_{h \to 0} \frac{f(x+h) - f(x)}{h}$.
    \item \textbf{Definite Integral:} The signed area under the curve of a function from $x=a$ to $x=b$, denoted as $\int_a^b f(x) \,dx$.
    \item \textbf{Partial Derivative:} The derivative of a multivariable function with respect to one variable, keeping the others constant.
\end{itemize}

\textbf{Theorems:}
\begin{itemize}
    \item \textbf{Fundamental Theorem of Calculus (FTC):} If $f$ is continuous on $[a,b]$ and $F$ is its antiderivative, then $\int_a^b f(x) \,dx = F(b) - F(a)$.
    \item \textbf{Green's Theorem:} Relates a line integral around a simple closed curve $C$ to a double integral over the plane region $D$ bounded by $C$: $\oint_C (L\,dx + M\,dy) = \iint_D \left(\frac{\partial M}{\partial x} - \frac{\partial L}{\partial y}\right) dA$.
\end{itemize}

\textbf{Examples:}
\begin{itemize}
    \item \textbf{Example:} Evaluate $\int_0^1 x^2 \,dx$. The antiderivative is $\frac{x^3}{3}$. Thus, the integral is $\frac{1^3}{3} - 0 = \frac{1}{3}$.
    \item \textbf{Example (Separable ODE):} Solve $\frac{dy}{dx} = y$. Separating variables gives $\frac{dy}{y} = dx$. Integrating both sides yields $\ln|y| = x + C$, or $y = Ce^x$.
\end{itemize}


\section{Linear Algebra}
\textbf{Prerequisite Knowledge:} Systems of linear equations, basic vector operations, and matrices.

\textbf{Definitions:}
\begin{itemize}
    \item \textbf{Vector Space:} A set of elements (vectors) that can be added together and multiplied by scalars, satisfying specific axioms like closure and associativity.
    \item \textbf{Eigenvalue and Eigenvector:} For a square matrix $A$, a non-zero vector $v$ is an eigenvector with eigenvalue $\lambda$ if $Av = \lambda v$.
    \item \textbf{Basis:} A set of linearly independent vectors that span a vector space.
\end{itemize}

\textbf{Theorems:}
\begin{itemize}
    \item \textbf{Rank-Nullity Theorem:} For any linear map $T: V \to W$, $\dim(\text{Im}(T)) + \dim(\text{Ker}(T)) = \dim(V)$.
    \item \textbf{Spectral Theorem:} Any real symmetric matrix can be diagonalized by an orthogonal matrix, meaning it has real eigenvalues and orthogonal eigenvectors.
\end{itemize}

\textbf{Examples:}
\begin{itemize}
    \item \textbf{Example:} Find the eigenvalues of $A = \begin{pmatrix} 2 & 0 \\ 0 & 3 \end{pmatrix}$. We solve $\det(A - \lambda I) = (2-\lambda)(3-\lambda) = 0$, giving eigenvalues $\lambda_1 = 2$, $\lambda_2 = 3$.
\end{itemize}


\section{Probability}
\textbf{Prerequisite Knowledge:} Basic set theory, combinatorics (permutations and combinations), and sequences.

\textbf{Definitions:}
\begin{itemize}
    \item \textbf{Random Variable:} A measurable function from a sample space to the real numbers.
    \item \textbf{Expectation (Mean):} The weighted average of all possible values a random variable can take. For a discrete variable, $E[X] = \sum_i x_i P(X = x_i)$.
    \item \textbf{Conditional Probability:} The probability of an event $A$ given that event $B$ has occurred, $P(A|B) = \frac{P(A \cap B)}{P(B)}$.
\end{itemize}

\textbf{Theorems:}
\begin{itemize}
    \item \textbf{Bayes' Theorem:} $P(A|B) = \frac{P(B|A)P(A)}{P(B)}$.
    \item \textbf{Central Limit Theorem (CLT):} The sum of a large number of independent, identically distributed random variables approaches a normal distribution, regardless of the original distribution.
\end{itemize}

\textbf{Examples:}
\begin{itemize}
    \item \textbf{Example:} Rolling a fair six-sided die. The sample space is $\{1,2,3,4,5,6\}$. The expectation is $E[X] = \frac{1+2+3+4+5+6}{6} = 3.5$.
\end{itemize}


\section{Operations Research}
\textbf{Prerequisite Knowledge:} Linear algebra, basic calculus, and systems of linear inequalities.

\textbf{Definitions:}
\begin{itemize}
    \item \textbf{Linear Programming (LP):} A method to achieve the best outcome (e.g., maximum profit or lowest cost) in a mathematical model whose requirements are represented by linear relationships.
    \item \textbf{Feasible Region:} The set of all possible points that satisfy the problem's constraints.
    \item \textbf{Zero-Sum Game:} A mathematical representation of a situation in which an advantage that is won by one of two sides is lost by the other.
\end{itemize}

\textbf{Theorems:}
\begin{itemize}
    \item \textbf{Strong Duality Theorem:} If either the primal or the dual linear programming problem has a finite optimal solution, then both have finite optimal solutions, and their optimal objective values are equal.
\end{itemize}

\textbf{Examples:}
\begin{itemize}
    \item \textbf{Example:} Maximize $Z = x + y$ subject to $x, y \geq 0$ and $x + 2y \leq 4$. The optimal solution occurs at a vertex of the feasible region, in this case $x=4, y=0$, giving $Z=4$.
\end{itemize}


\section{Finite Difference Methods for PDEs}
\textbf{Prerequisite Knowledge:} Partial differential equations, Taylor series expansion.

\textbf{Definitions:}
\begin{itemize}
    \item \textbf{Finite Difference:} A mathematical expression of the form $f(x+h) - f(x)$. Divided by $h$, it approximates the derivative.
    \item \textbf{Consistency:} A numerical scheme is consistent if the local truncation error goes to zero as the grid spacing goes to zero.
    \item \textbf{Stability:} A numerical scheme is stable if errors do not grow exponentially as the computation proceeds.
\end{itemize}

\textbf{Theorems:}
\begin{itemize}
    \item \textbf{Lax Equivalence Theorem:} For a consistent finite difference method for a well-posed linear initial value problem, the method is convergent if and only if it is stable.
\end{itemize}

\textbf{Examples:}
\begin{itemize}
    \item \textbf{Example:} The Forward Time Centered Space (FTCS) scheme for the 1D heat equation $u_t = \alpha u_{xx}$ is $\frac{u_i^{n+1} - u_i^n}{\Delta t} = \alpha \frac{u_{i+1}^n - 2u_i^n + u_{i-1}^n}{(\Delta x)^2}$.
\end{itemize}


\section{Finite Element Methods (FEM)}
\textbf{Prerequisite Knowledge:} Vector spaces, calculus of variations, linear boundary value problems.

\textbf{Definitions:}
\begin{itemize}
    \item \textbf{Weak Formulation:} Rewriting a differential equation in an integral form that requires less smoothness from the solution.
    \item \textbf{Galerkin Method:} A class of methods for converting a continuous operator problem to a discrete problem by restricting the weak formulation to a finite-dimensional subspace.
\end{itemize}

\textbf{Theorems:}
\begin{itemize}
    \item \textbf{C\'ea's Lemma:} Establishes that the Galerkin approximation is the best approximation in the finite-dimensional subspace, up to a constant factor.
\end{itemize}

\textbf{Examples:}
\begin{itemize}
    \item \textbf{Example:} To solve $-u'' = f$ with $u(0)=u(1)=0$, the weak form is: find $u$ such that $\int_0^1 u' v' \,dx = \int_0^1 f v \,dx$ for all test functions $v$.
\end{itemize}


\section{Numerical Linear Algebra}
\textbf{Prerequisite Knowledge:} Linear algebra, matrix norms, and orthogonality.

\textbf{Definitions:}
\begin{itemize}
    \item \textbf{Singular Value Decomposition (SVD):} A factorization of a real or complex matrix $A$ into $U \Sigma V^T$, where $U$ and $V$ are orthogonal and $\Sigma$ is diagonal.
    \item \textbf{Condition Number:} A measure of how sensitive the answer is to small changes in the input data.
\end{itemize}

\textbf{Theorems:}
\begin{itemize}
    \item \textbf{Eckart-Young Theorem:} The best rank-$k$ approximation to a matrix $A$ under the Frobenius norm is given by the truncated SVD of $A$.
    \item \textbf{Gershgorin Circle Theorem:} Every eigenvalue of a matrix lies within at least one of the Gershgorin discs constructed from the matrix entries.
\end{itemize}

\textbf{Examples:}
\begin{itemize}
    \item \textbf{Example:} The Power Method repeatedly multiplies a vector by a matrix $A$ and normalizes it to find the dominant eigenvalue and its corresponding eigenvector.
\end{itemize}


\section{Numerical Methods for ODEs}
\textbf{Prerequisite Knowledge:} Ordinary differential equations, Taylor series, and basic numerical analysis.

\textbf{Definitions:}
\begin{itemize}
    \item \textbf{Explicit Method:} A numerical method where the new state $y_{n+1}$ is computed directly from known previous states.
    \item \textbf{Implicit Method:} A numerical method where the new state $y_{n+1}$ is defined by an equation involving $y_{n+1}$ itself.
    \item \textbf{A-stability:} A stability property essential for solving stiff differential equations.
\end{itemize}

\textbf{Theorems:}
\begin{itemize}
    \item \textbf{Dahlquist's Second Barrier:} No explicit linear multistep method is A-stable. The highest order of an A-stable implicit linear multistep method is 2.
\end{itemize}

\textbf{Examples:}
\begin{itemize}
    \item \textbf{Example:} Forward Euler for $y' = f(t,y)$ is given by $y_{n+1} = y_n + \Delta t f(t_n, y_n)$.
\end{itemize}


\section{Optimization and Convex Analysis}
\textbf{Prerequisite Knowledge:} Multivariable calculus, linear algebra, and real analysis.

\textbf{Definitions:}
\begin{itemize}
    \item \textbf{Convex Function:} A function where the line segment between any two points on its graph lies above or on the graph.
    \item \textbf{Lagrangian:} A function that combines the objective function and the constraints of an optimization problem using Lagrange multipliers.
\end{itemize}

\textbf{Theorems:}
\begin{itemize}
    \item \textbf{Karush-Kuhn-Tucker (KKT) Conditions:} Necessary (and under convexity, sufficient) conditions for a solution to be optimal in nonlinear programming.
    \item \textbf{Slater's Condition:} A sufficient condition for strong duality to hold in a convex optimization problem.
\end{itemize}

\textbf{Examples:}
\begin{itemize}
    \item \textbf{Example:} Minimize $f(x) = x^2$ subject to $x \ge 1$. The KKT conditions quickly show the minimum is at the boundary $x=1$.
\end{itemize}


\section{Calculus of Variations and Euler-Lagrange}
\textbf{Prerequisite Knowledge:} Advanced calculus, integration by parts, and differential equations.

\textbf{Definitions:}
\begin{itemize}
    \item \textbf{Functional:} A mapping from a vector space of functions to the real numbers (e.g., an integral involving an unknown function).
    \item \textbf{First Variation:} The generalization of the directional derivative to functionals.
\end{itemize}

\textbf{Theorems:}
\begin{itemize}
    \item \textbf{Euler-Lagrange Equation:} The fundamental equation of the calculus of variations. If $J[y] = \int_a^b L(x, y, y')\,dx$ is extremized by $y$, then $\frac{\partial L}{\partial y} - \frac{d}{dx}\left(\frac{\partial L}{\partial y'}\right) = 0$.
\end{itemize}

\textbf{Examples:}
\begin{itemize}
    \item \textbf{Example:} The shortest path between two points in a plane. The functional is length $L = \int \sqrt{1 + (y')^2} dx$. The Euler-Lagrange equation implies $y'' = 0$, meaning the path is a straight line.
\end{itemize}


\section{Orthogonal Polynomials and Spectral Methods}
\textbf{Prerequisite Knowledge:} Inner product spaces, integration techniques, and differential equations.

\textbf{Definitions:}
\begin{itemize}
    \item \textbf{Orthogonal Polynomials:} A family of polynomials such that any two different polynomials in the sequence are orthogonal to each other under an inner product.
    \item \textbf{Gaussian Quadrature:} A numerical integration rule that yields exact results for polynomials up to a certain degree by carefully choosing evaluation points and weights.
\end{itemize}

\textbf{Theorems:}
\begin{itemize}
    \item \textbf{Exactness of Gaussian Quadrature:} An $n$-point Gaussian quadrature rule integrates polynomials of degree up to $2n - 1$ exactly.
\end{itemize}

\textbf{Examples:}
\begin{itemize}
    \item \textbf{Example:} Legendre polynomials $P_n(x)$ are orthogonal on $[-1,1]$. The first few are $P_0(x) = 1$ and $P_1(x) = x$.
\end{itemize}


\section{Dynamical Systems and Biological Modeling}
\textbf{Prerequisite Knowledge:} Systems of ODEs, linear algebra (eigenvalues), and basic biology concepts for modeling.

\textbf{Definitions:}
\begin{itemize}
    \item \textbf{Fixed Point (Equilibrium):} A point $x^*$ in the state space where the system does not change, i.e., $f(x^*) = 0$ for $\frac{dx}{dt} = f(x)$.
    \item \textbf{Bifurcation:} A sudden qualitative change in the behavior of a dynamical system as a parameter is varied.
\end{itemize}

\textbf{Theorems:}
\begin{itemize}
    \item \textbf{Linearization Theorem (Hartman-Grobman):} Near a hyperbolic equilibrium point, the nonlinear system behaves qualitatively like its linearization.
\end{itemize}

\textbf{Examples:}
\begin{itemize}
    \item \textbf{Example:} Logistic growth model $\frac{dx}{dt} = r x(1 - x/K)$. The fixed points are $x=0$ (unstable) and $x=K$ (stable carrying capacity).
\end{itemize}


\section{Graph Theory}
\textbf{Prerequisite Knowledge:} Basic set theory and discrete mathematics.

\textbf{Definitions:}
\begin{itemize}
    \item \textbf{Graph:} A set of vertices connected by edges.
    \item \textbf{Eulerian Circuit:} A trail in a finite graph that visits every edge exactly once and starts and ends on the same vertex.
    \item \textbf{Planar Graph:} A graph that can be drawn on the plane in such a way that its edges intersect only at their endpoints.
\end{itemize}

\textbf{Theorems:}
\begin{itemize}
    \item \textbf{Euler's Formula:} For any connected planar graph with $V$ vertices, $E$ edges, and $F$ faces, $V - E + F = 2$.
    \item \textbf{Dirac's Theorem:} If a simple graph has $n \ge 3$ vertices and every vertex has degree at least $n/2$, then the graph is Hamiltonian.
\end{itemize}

\textbf{Examples:}
\begin{itemize}
    \item \textbf{Example:} The complete graph on 4 vertices ($K_4$) is planar, but $K_5$ is not, by Kuratowski's theorem.
\end{itemize}


\section{Number Theory}
\textbf{Prerequisite Knowledge:} Basic arithmetic, algebra, and induction.

\textbf{Definitions:}
\begin{itemize}
    \item \textbf{Prime Number:} A natural number greater than 1 that has no positive divisors other than 1 and itself.
    \item \textbf{Congruence:} $a \equiv b \pmod m$ means that $m$ divides the difference $a - b$.
\end{itemize}

\textbf{Theorems:}
\begin{itemize}
    \item \textbf{Fermat's Little Theorem:} If $p$ is a prime number, then for any integer $a$, the number $a^p - a$ is an integer multiple of $p$.
    \item \textbf{Chinese Remainder Theorem:} Provides a unique solution to simultaneous linear congruences with coprime moduli.
\end{itemize}

\textbf{Examples:}
\begin{itemize}
    \item \textbf{Example:} Solve $x \equiv 2 \pmod 3$ and $x \equiv 3 \pmod 5$. A solution is $x = 8$, since $8 = 2+2\cdot3$ and $8 = 3+1\cdot5$.
\end{itemize}


\section{Learning Theory and VC Dimension}
\textbf{Prerequisite Knowledge:} Probability theory, set theory, and basics of machine learning.

\textbf{Definitions:}
\begin{itemize}
    \item \textbf{PAC Learning:} Probably Approximately Correct learning, a framework for mathematical analysis of machine learning.
    \item \textbf{VC Dimension (Vapnik-Chervonenkis Dimension):} A measure of the capacity (complexity, expressive power, richness, or flexibility) of a space of functions that can be learned by a statistical classification algorithm.
\end{itemize}

\textbf{Theorems:}
\begin{itemize}
    \item \textbf{Sauer's Lemma:} Binds the growth function by a polynomial in $N$ whose degree is the VC dimension $d$.
    \item \textbf{Fundamental Theorem of PAC Learning:} A hypothesis class is PAC learnable if and only if it has a finite VC dimension.
\end{itemize}

\textbf{Examples:}
\begin{itemize}
    \item \textbf{Example:} The VC dimension of linear classifiers (lines) in 2D space is 3, because they can shatter any 3 non-collinear points but no set of 4 points.
\end{itemize}
